% TODO make hw stylesheet
\documentclass[11pt]{scrartcl}
\usepackage[a4paper, total={6in, 8in}]{geometry}
\usepackage[usenames,dvipsnames,svgnames]{xcolor}
\usepackage[shortlabels]{enumitem}
\usepackage[framemethod=TikZ]{mdframed}
\usepackage{amsmath,amssymb,amsthm}
\usepackage[colorlinks]{hyperref}
\usepackage{multicol}
\usepackage{tabularx}

\hypersetup{
	colorlinks=true,
	linkcolor=blue,
	filecolor=magenta,      
	urlcolor=magenta,
	pdftitle={MAT 362 HW}
}

\usepackage{microtype}
\usepackage{mathtools}
\usepackage[headsepline]{scrlayer-scrpage}
\usepackage{thmtools}
\usepackage[textsize=footnotesize]{todonotes}

\mdfdefinestyle{mdbluebox}{roundcorner=10pt,innerbottommargin=9pt,
	linecolor=blue,backgroundcolor=TealBlue!5,}
\declaretheoremstyle[headfont=\sffamily\bfseries\color{MidnightBlue},
mdframed={style=mdbluebox},]{thmbluebox}

\mdfdefinestyle{mdbloobox}{frametitlefont=\bfseries,innerbottommargin=8pt,
	nobreak=true,backgroundcolor=TealBlue!5,linecolor=blue,}
\declaretheoremstyle[headfont=\bfseries\color{MidnightBlue},
mdframed={style=mdbloobox},headpunct={\\[3pt]},postheadspace=0pt,]{thmbloobox}

\mdfdefinestyle{mdredbox}{frametitlefont=\bfseries,innerbottommargin=8pt,
	nobreak=true,backgroundcolor=Salmon!5,linecolor=RawSienna,}
\declaretheoremstyle[headfont=\bfseries\color{RawSienna},
mdframed={style=mdredbox},headpunct={\\[3pt]},postheadspace=0pt,]{thmredbox}

\mdfdefinestyle{mdgreenbox}{linecolor=ForestGreen,backgroundcolor=ForestGreen!5,
	linewidth=2pt,rightline=false,leftline=true,topline=false,bottomline=false,}
\declaretheoremstyle[headfont=\bfseries\sffamily\color{ForestGreen!70!black},
mdframed={style=mdgreenbox},headpunct={ --- },]{thmgreenbox}

\mdfdefinestyle{mdorangebox}{linecolor=RedOrange,backgroundcolor=RedOrange!5,
	linewidth=2pt,rightline=false,leftline=true,topline=false,bottomline=false,}
\declaretheoremstyle[headfont=\bfseries\sffamily\color{RedOrange!70!black},
mdframed={style=mdorangebox},headpunct={ --- },]{thmorangebox}

\mdfdefinestyle{mdblackbox}{linecolor=black,backgroundcolor=RedViolet!5!gray!5,
	linewidth=3pt,rightline=false,leftline=true,topline=false,bottomline=false,}
\declaretheoremstyle[mdframed={style=mdblackbox}]{thmblackbox}

\declaretheoremstyle[spaceabove=6pt,spacebelow=6pt]{basehead}

\declaretheorem[style=basehead,name=Problem]{problem}
\declaretheorem[style=basehead,name=Problem,numbered=no]{problem*}
\declaretheorem[style=thmbloobox,name=Setup]{setup} % for config geo
\declaretheorem[style=thmredbox,name=Example,sibling=problem]{example}
\declaretheorem[style=thmredbox,name=$\bigstar$ Problem,sibling=problem]{niceprob}
\declaretheorem[style=thmbluebox,name=Theorem,sibling=problem]{theorem}
\declaretheorem[style=thmbluebox,name=Corollary,sibling=theorem]{corollary}
\declaretheorem[style=thmbluebox,name=Proposition,sibling=theorem]{prop}
\declaretheorem[style=thmbluebox,name=Lemma,numberwithin=problem]{lemma}
\declaretheorem[style=thmbluebox,name=Theorem,numbered=no]{theorem*}
\declaretheorem[style=thmbluebox,name=Lemma,numbered=no]{lemma*}
\declaretheorem[style=thmgreenbox,name=Claim,numberwithin=problem]{claim}
\declaretheorem[style=thmgreenbox,name=Claim,numbered=no]{claim*}
\declaretheorem[style=thmblackbox,name=Remark,numberwithin=problem]{remark}
\declaretheorem[style=thmblackbox,name=Remark,numbered=no]{remark*}
\declaretheorem[style=thmgreenbox,name=Definition,sibling=theorem]{definition}
\declaretheorem[style=thmgreenbox,name=Definition,numbered=no]{definition*}
\declaretheorem[style=thmorangebox,name=Warning,sibling=theorem]{warning}
\declaretheorem[style=thmorangebox,name=Warning,numbered=no]{warning*}
\declaretheorem[style=thmorangebox,name=Note,numbered=no]{note}
\declaretheorem[style=thmblackbox,name=Exercise,sibling=problem]{exercise}
\declaretheorem[style=thmblackbox,name=Exercise,numbered=no]{exercise*}
\declaretheorem[style=thmblackbox,name=Question,sibling=problem]{question}
\declaretheorem[style=thmblackbox,name=Question,numbered=no]{question*}

\addtolength{\textheight}{3.14cm}
\providecommand{\ol}{\overline}
\providecommand{\ul}{\underline}
\providecommand{\eps}{\varepsilon}
\providecommand{\half}{\frac{1}{2}}
\providecommand{\dang}{\measuredangle} %% Directed angle
\providecommand{\CC}{\mathbb C}
\providecommand{\FF}{\mathbb F}
\providecommand{\NN}{\mathbb N}
\providecommand{\QQ}{\mathbb Q}
\providecommand{\RR}{\mathbb R}
\providecommand{\ZZ}{\mathbb Z}
\providecommand{\dg}{^\circ}
\providecommand{\ii}{\item}
\providecommand{\setdiff}{\smallsetminus}
\providecommand{\alert}[1]{{\sffamily\textbf{\textcolor{blue}{#1}}}}
\providecommand{\inv}{^{-1}}
\providecommand{\mb}[1]{\mathbf{#1}}
\providecommand{\dif}{\;d}
\newcommand*{\horzbar}{\rule[.5ex]{2.5ex}{0.5pt}}

\reversemarginpar

\newenvironment{soln}{\begin{proof}[Solution]}{\end{proof}}
\newenvironment{walkthrough}{\noindent \textbf{Walkthrough.}}{\medskip}
\newlist{walk}{enumerate}{3}
\setlist[walk]{label=\bfseries (\alph*)}

\providecommand{\pow}{\operatorname{Pow}}
\providecommand{\vocab}[1]{{\textbf{\textcolor{ForestGreen}{#1}}}}
\providecommand{\lcm}{\operatorname{lcm}}
\providecommand{\abs}[1]{\left|#1\right|}

\usepackage{asymptote}
\begin{asydef}
	defaultpen(fontsize(10pt));
	unitsize(1cm);
	size(8cm); // set a reasonable default
	usepackage("amsmath");
	usepackage("amssymb");
	settings.tex="pdflatex";
	settings.outformat="pdf";
	// Replacement for olympiad+cse5 which is not standard
	import geometry;
	// recalibrate fill and filldraw for conics
	void filldraw(picture pic = currentpicture, conic g, pen fillpen=defaultpen, pen drawpen=defaultpen)
	{ filldraw(pic, (path) g, fillpen, drawpen); }
	void fill(picture pic = currentpicture, conic g, pen p=defaultpen)
	{ filldraw(pic, (path) g, p); }
	// some geometry
	pair foot(pair P, pair A, pair B) { return foot(triangle(A,B,P).VC); }
	pair orthocenter(pair A, pair B, pair C) { return orthocentercenter(A,B,C); }
	pair centroid(pair A, pair B, pair C) { return (A+B+C)/3; }
	// cse5 abbreviations
	path CP(pair P, pair A) { return circle(P, abs(A-P)); }
	path CR(pair P, real r) { return circle(P, r); }
	pair IP(path p, path q) { return intersectionpoints(p,q)[0]; }
	pair OP(path p, path q) { return intersectionpoints(p,q)[1]; }
	path Line(pair A, pair B, real a=0.6, real b=a) { return (a*(A-B)+A)--(b*(B-A)+B); }
	// cse5 more useful functions
	picture CC() {
		picture p=rotate(0)*currentpicture;
		currentpicture.erase();
		return p;
	}
	pair MP(Label s, pair A, pair B = plain.S, pen p = defaultpen) {
		Label L = s;
		L.s = "$"+s.s+"$";
		label(L, A, B, p);
		return A;
	}
	pair Drawing(Label s = "", pair A, pair B = plain.S, pen p = defaultpen) {
		dot(MP(s, A, B, p), p);
		return A;
	}
	path Drawing(path g, pen p = defaultpen, arrowbar ar = None) {
		draw(g, p, ar);
		return g;
	}
\end{asydef}

\usepackage{mathrsfs}

\ihead{\footnotesize MAT 362 HW02}
\ohead{\footnotesize Last compiled \today}

\title{\vspace{-2em}MAT 362 HW02}
\author{\large Aditya Pahuja}
\date{\large Due 02/13}
\begin{document}
\maketitle

\begin{problem*}1.5.1.
	Given the parametrized curve
	\[\alpha(s) = \left(a\cos\frac sc, a\sin\frac sc, b\frac sc\right),\quad
	s\in\RR,\]
	where $c^2 = a^2 + b^2$,
	\begin{enumerate}[(a)]
		\ii Show that the parameter $s$ is the arc length.
		\ii Determine the curvature and torsion of $\alpha$.
		\ii Determine the osculating plane of $\alpha$.
		\ii Show that the lines containing $n(s)$ and
		passing through $\alpha(s)$ meet the $z$-axis under a constant angle
		equal to $\pi/2$.
		\ii Show that the tangent lines to $\alpha$ make
		a constant angle with the $z$-axis.
	\end{enumerate}
\end{problem*}

\begin{soln}
	(a) It suffices to show that $|\alpha'(s)| = 1$. To see this, we compute
	\[|\alpha'(s)| = \sqrt{\left(-\frac ac\sin\frac sc\right)^2 +
	\left(\frac ac\cos\frac sc\right)^2 + \left(\frac bc\right)^2}
	= \sqrt{\frac{a^2 + b^2}{c^2}} = 1.\]
	
	(b) The curvature of $\alpha$ is equal to $|\alpha''(s)|$,
	which in this case is
	\[\abs{\left(-\frac{a}{c^2}\cos\frac sc, -\frac{a}{c^2}\sin\frac sc,
		0\right)} = \boxed{\frac{a}{c^2}}.\]
	The normal vector $n(s)$ is then $(-\cos\frac sc,-\sin\frac sc,0)$
	and we know that $t(s) = (-\frac ac\sin\frac sc, \frac ac\cos\frac sc,
	\frac bc)$, so their vector product is
	\[b(s)= t(s)\wedge n(s) = \left(\frac bc\sin\frac sc, -\frac bc\cos\frac sc,
	\frac ac\right),\]
	so the torsion of $\alpha$ is
	\[|b'(s)| = \abs{\left(\frac b{c^2}\cos\frac sc,
		\frac b{c^2}\cos\frac sc,0\right)}
	= \boxed{\frac b{c^2}}.\]
	
	(c) The osculating plane of $\alpha$ is the plane through $\alpha(s)$
	determined by $t(s)$ and $n(s)$, i.e. orthogonal to $b(s)$.
	This means it is the set of points $x$ such that
	\[(x - \alpha(s))\cdot b(s) = 0.\]
	Expanding this yields the plane
	\[\boxed{x_1\cdot b\sin\frac sc - x_2\cdot b\cos\frac sc + ax_3 = abs}.\]
	
	(d) We computed above that $n(s) = (-\cos\frac sc, -\sin\frac sc, 0)$,
	so the line in the direction of $n(s)$ passing through $\alpha(s)$ is
	\[\ell(x) = \alpha(s) + x\cdot n(s)\]
	for $x\in\RR$. This line meets the $z$-axis at $x = a$, which is the point
	\[\ell(a) = \left(0, 0, \frac{bs}c\right).\]
	It is clear that the line makes a right angle with the $z$-axis
	because $n(s)\cdot (0,0,1) = 0$ means that $n(s)$
	and the unit vector in the direction of the $z$-axis are orthogonal.
	
	(e) The tangent lines to $\alpha$ make an angle of
	\[\arccos(t(s)\cdot (0,0,1)) =
	\arccos\left(
	\left(-\frac ac\sin\frac sc, \frac ac\cos\frac sc, \frac bc\right)
	\cdot(0,0,1)\right) = \arccos\frac bc\]
	with the $z$-axis, which is independent of our only parameter $s$.
\end{soln}

\begin{problem*}1.5.6.
	A \vocab{translation} by a vector $v\in\RR^3$ is the map
	$A\colon\RR^3\to\RR^3$ that is given by $A(p) = p + v$.
	A linear map $\rho\colon\RR^3\to\RR^3$ is an
	\vocab{orthogonal transformation} if $\rho u\cdot\rho v = u\cdot v$
	for all $u,v\in\RR^3$. A \vocab{rigid motion} in $\RR^3$ is
	the result of composing a translation with an orthogonal transformation
	with positive determinant.
	\begin{enumerate}[(a)]
		\ii Demonstrate that the norm of a vector
		and the angle $\theta$ between two vectors, $0\le \theta\le\pi$,
		are invariant under orthogonal transformations with
		positive determinant.
		\ii Show that the vector product of two vectors is invariant under
		orthogonal transformations with positive determinant.
		is the assertion still true if we drop the condition
		on the determinant?
		\ii Show that the arc length, the curvature, and the torsion
		of a parametrized curve are (whenever defined)
		invariant under rigid motions.
	\end{enumerate}
\end{problem*}

\begin{soln}
	(a) Let $\rho$ be an orthogonal transformation and let $u\in\RR^3$. Then
	\[\|u\|^2 = u\cdot u = \rho u\cdot\rho u = \|\rho u\|^2\]
	implies norm is preserved under orthogonal transformations.
	Moreover, if $\theta$ is the angle between two vectors $u,v\in\RR^3$,
	\[\cos\theta = \frac{u\cdot v}{\|u\|\|v\|}
	= \frac{\rho u\cdot \rho v}{\|\rho u\|\|\rho v\|} = \cos\theta'\]
	where $\theta'$ is the angle between $\rho u$ and $\rho v$
	(with $0\le\theta, \theta' \le \pi$), which implies that
	$\theta = \theta'$ because $\cos$ is injective on $[0,\pi]$,
	i.e. angle is preserved by $\rho$.
	
	(b) Let $u,v\in\RR^3$ be two vectors. The above problem shows
	that $\rho u$, $\rho v$, and $\rho (u\wedge v)$
	have the same norms as $u$, $v$, $u\wedge v$, and moreover,
	$\rho(u\wedge v)$ is orthogonal to $\rho u$ and $\rho v$.
	This narrows down $\rho(u\wedge v)$ to two possible vectors:
	either $\rho u\wedge \rho v$ or $-\rho u\wedge\rho v$.
	This is where we can use the fact that $\rho$ has positive determinant:
	since $\rho$ preserves orientation,
	$\rho u$, $\rho v$, and $\rho (u\wedge v)$ have the same orientation
	as $u$, $v$, and $u\wedge v$,
	which confirms that $\rho(u\wedge v) = \rho u\wedge\rho v$.
	
	Indeed, if $\rho$ is allowed to have negative orientation, then
	consider the transformation given by the matrix
	\[\rho = \begin{pmatrix}
		1&0&0\\0&1&0\\0&0&-1
	\end{pmatrix}\]
	in the standard basis, i.e. reflection through the $xy$-plane.
	This is an orthogonal transformation because for any two $u,v\in\RR^3$,
	\[\rho u\cdot\rho v = (u_1,u_2,-u_3)\cdot(v_1,v_2,-v_3) =
	u_1v_1 + u_2v_2 + u_3v_3 =
	u\cdot v.\]
	However, if $e_1$, $e_2$, $e_3$ are the standard basis vectors, then
	\[e_3 = e_1\wedge e_2 = \rho e_1 \wedge \rho e_2 = -\rho e_3\]
	because $\rho e_1 = e_1$ and $\rho e_2 = e_2$,
	so $\rho$ doesn't preserve vector products.
	
	(c) Let $\alpha(s)$ be an arc length parametrized curve.
	Let $\rho$ be an orthogonal transformation with matrix
	\[\rho = \begin{pmatrix}
		\horzbar&v_1&\horzbar\\
		\horzbar&v_2&\horzbar\\
		\horzbar&v_3&\horzbar
	\end{pmatrix}\]
	where $v_1$, $v_2$, $v_3$ are row vectors.
	To start with we show that tangent vectors are invariant under $\rho$:
	\[(\rho\alpha)'(s) = \begin{pmatrix}
		v_1\cdot\alpha(s)\\v_2\cdot\alpha(s)\\v_3\cdot\alpha(s)
	\end{pmatrix}' =\begin{pmatrix}
		v_1\cdot\alpha'(s)\\v_2\cdot\alpha'(s)\\v_3\cdot\alpha'(s)
	\end{pmatrix} = \rho\cdot\alpha'(s).\]
	In particular, the arc length of the curve $\rho\alpha(s)$
	from $s=a$ to $s=b$ is
	\[\int_a^b |\rho\alpha'(s)|\dif s = \int_a^b |\alpha'(s)|\dif s,\]
	recalling that $\rho$ preserves norm, so arc length is preserved by $\rho$.
	Then, the curvature is
	\[|(\rho\alpha)''(s)| = |\rho\cdot\alpha''(s)| =
	|\alpha''(s)| = \kappa(s),\]
	so it is also preserved by $\rho$. For the torsion, we can check first that
	the binormal of $\alpha$ at point $s$ is mapped to the binormal
	of $\rho\alpha$ at $s$: clearly $t(s) = \alpha'(s)$
	and $n(s) = \alpha''(s)/\kappa(s)$ are preserved by $\rho$,
	and part (b) says that their vector product is also preserved by $\rho$.
	Then we have that the torsion of $\rho\alpha$ is the norm of
	$(\rho b)'(s)$, which is again just $|b'(s)| = \tau(s)$,
	so finally torsion is preserved by $\rho$.
\end{soln}

\begin{problem*}1.5.8(a).
	The trace of the parametrized curve (arbitrary parameter)
	\[\alpha(t) = (t, \cosh t),\quad t\in\RR\]
	is called the \vocab{catenary}. Show that the signed curvature
	of the catenary is $\kappa(t) = \frac{1}{\cosh^2t}$.
\end{problem*}

\begin{soln}
	Recall that $\cosh t = \half(e^t + e^{-t})$ and
	$\sinh t = \half(e^t - e^{-t})$. We can compute
	\[\alpha'(t) = (1, \sinh t),\quad |\alpha'(t)| = \sqrt{1 + \sinh^2 t}
	= \cosh t\]
	(note that we do not need to say $|\cosh t|$ because $e^t$ and $e^{-t}$
	are always positive).Therefore the normalized tangent vector is
	\[T(t) = \left(\frac{1}{\cosh t}, \tanh t\right).\]
	We then compute
	\[\frac{d}{ds} T(t) = \frac{dt}{ds}\cdot T'(t).\]
	We know that
	\[\frac{dt}{ds} = \frac{1}{ds/dt} = \frac{1}{|\alpha'(t)|},\]
	so the derivative of $T$ with respect to $s$ is
	\[\frac{T'(t)}{|\alpha'(t)|} = \frac{1}{\cosh t}\cdot
	\left(-\frac{\sinh t}{\cosh^2t}, 1-\frac{\sin^2t}{\cosh^2t}\right) =
	\frac{1}{\cosh^2 t}\left(-\tanh t, \frac{1}{\cosh t}\right),\]
	noting that $\cosh^2 t - \sinh^2 t = 1$.
	On the other hand, the unit normal vector $n$ such that
	$\{T, n\}$ has the same orientation as $\{e_1, e_2\}$ is the
	$\frac\pi2$-counterclockwise rotation of $T$, which is
	\[\left(-\tanh t, \frac{1}{\cosh t}\right).\]
	By definition, the signed curvature is the scalar $\kappa$ such that
	$T'(s) = \kappa n$, so we find that $\kappa = \frac{1}{\cosh^2t}$
	as desired.
\end{soln}

\begin{problem*} 1.5.10.
	Consider the map
	\[\alpha(t) = \begin{cases}
		(t, 0, e^{-1/t^2}), & t>0\\
		(t, e^{-1/t^2}, 0), & t<0\\
		(0, 0, 0) & t=0
	\end{cases}\]
	\begin{enumerate}[(a)]
		\ii Prove that $\alpha$ is a differentiable curve.
		\ii Prove that $\alpha$ is regular for all $t$ and that
		the curvature $\kappa(t)\ne 0$ for $t\ne 0$, $t\ne\pm\sqrt{2/3}$,
		and $\kappa(0) = 0$.
		\ii Show that the limit of the osculating planes as $t\to 0$,
		$t > 0$ is the plane $y=0$ but that the limit of
		the osculating planes as $t\to 0$, $t<0$,
		is the plane $z=0$.
		\ii Show that $\tau$ can be defined so that $\tau\equiv 0$,
		even though $\alpha$ is not a plane curve.
	\end{enumerate}
\end{problem*}

\begin{soln}
	(a) Differentiability away from zero is trivial, so we just check
	that $\alpha'(0)$ exists. Specifically, we need to check that the functions
	\[y(t) = \begin{cases}
		e^{-1/t^2} & t < 0\\
		0 & t\ge 0
	\end{cases}\qquad z(t) = \begin{cases}
		e^{-1/t^2} & t > 0\\
		0 & t\le 0
	\end{cases}\]
	are differentiable at zero. Notice that $y(t) = z(-t)$,
	so we will verify differentiability for $y$ only,
	since differentiability at $0$
	for one function implies the same for the other.
	This amounts to computing
	\[\lim_{h\to 0} \frac{y(h) - y(0)}{h}.\]
	From the left, this limit is clearly $0$, so we need to show that
	\[\lim_{h\to 0^+} \frac{e^{-1/h^2}}{h} = 0.\]
	We have
	\[0\le\lim_{h\to 0^+} \frac{e^{-1/h^2}}{h}=
	\lim_{x\to\infty} \frac{x}{e^{x^2}} \le
	\lim_{x\to\infty} \frac{x}{x^2} = 0\]
	where the latter inequality comes from $e^r > r^2$ for large enough $r$,
	so this confirms that $y'(0) = z'(0) = 0$.
	
	(b) $\alpha$ is regular for all $t$ because the $x$-component
	of $\alpha'$ is always $1$, so $|\alpha'| \ge 1$ for all $t$.
	To compute the curvature we use the following fact,
	noting that $\alpha$ is essentially a plane curve
	when restricting to $t > 0$ or $t < 0$.
	\begin{claim*}[Problem 1.5.12(d)]
		If $\alpha\colon I\to\RR^2$ is a plane curve
		$\alpha(t) = (x(t), y(t))$, the signed curvature of $\alpha$ at $t$ is
		\[\kappa(t) = \frac{x'y'' - x''y'}{((x')^2 + (y')^2)^{3/2}}.\]
	\end{claim*}
	\begin{proof}
		To start with we can calculate the unit tangent
		\[T(t) = \left(\frac{x'}{\sqrt{(x')^2 + (y')^2}},
		\frac{y'}{\sqrt{(x')^2 + (y')^2}}\right).\]
		Then, the derivative of the unit tangent with respect to $t$ is
		\[\left(\frac{-y'(x'y'' - x''y')}{((x')^2 + (y')^2)^{3/2}},
		\frac{x'(x'y'' - x''y')}{((x')^2 + (y')^2)^{3/2}}\right),\]
		which can be verified by some menial computation.
		To get the derivative with respect to $s$, we just multiply this by
		$dt/ds = \frac1{|\alpha'(t)|}$. This means that,
		if $n$ is the normal vector obtained by a $\frac\pi2$-counterclockwise
		rotation of $T$, then the signed curvature is the scalar $\kappa$
		such that $T = \kappa n$: in this case
		\[n = \left(\frac{-y'}{|\alpha'|}, \frac{x'}{|\alpha'|}\right),\]
		and we know that
		\[\frac{dT}{ds} = \frac{1}{|\alpha'|}\cdot
		\frac{x'y'' - x''y'}{((x')^2 + (y')^2)^{3/2}}\cdot(-y',x') =
		\frac{x'y'' - x''y'}{((x')^2 + (y')^2)^{3/2}}\cdot
		\frac{(-y', x')}{|\alpha'|},\]
		so we can see that the curvature is indeed
		the expression in the claim.
	\end{proof}
	
	We focus on $t\ge 0$, since $t < 0$ and $t > 0$
	are symmetrical and so the computations for a given value of $t$
	end up being the same as for $-t$.
	For $t > 0$, the curvature is equal to
	\[\frac{|x'z'' - x''z'|}{|\alpha'|^3}
	= \frac{|z''|}{|\alpha'|^3},\]
	so we need to find $t > 0$ for which $z'' = 0$ in order to check
	when the curvature is $0$. We have
	\[z''(t) = \frac{-6}{t^4}e^{-1/t^2} + \frac{4}{t^6} e^{-1/t^2},\]
	which is zero if and only if
	\[6t^2 = 4,\]
	so for $t > 0$ the only point of zero curvature is $t = \sqrt{2/3}$.
	Symmetrically, for $t<0$, the only point of zero curvature
	is $t = -\sqrt{2/3}$.
	
	Finally, for $t = 0$, we can tediously compute that
	the normal vector is the zero vector, which is enough to imply
	that the curvature is $0$.
	
	(c) For $t > 0$, the $y$-coordinate is always $0$,
	so the curve as well as its tangent and normal vectors
	live in the $y=0$ plane,
	hence the osculating plane is the $y=0$ plane as well;
	similarly, the curve and its tangent and normal vectors
	all lie in the $z=0$ plane for $t<0$, so
	the osculating plane there is $z=0$.
	The corresponding limits follow immediately.
	
	(d) The torsions for $t > 0$ and for $t < 0$ are both zero
	because $\alpha$ lies in a single plane for each of $t > 0$ and $t < 0$.
	Moreover we can continuously extend $\tau$ to be $0$ at $t=0$,
	so we can reasonably have $\tau\equiv 0$ despite $\alpha$
	not being a planar curve.
\end{soln}

\begin{problem*}1.5.11.
	One often gives a plane curve in polar coordinates by $\rho = \rho(\theta)$,
	$a\le\theta\le b$.
	\begin{enumerate}[(a)]
		\ii Show that the arc length is
		\[\int_a^b \sqrt{\rho^2 + (\rho')^2}\dif\theta,\]
		where the prime denotes the derivative relative to $\theta$.
		\ii Show that the curvature is
		\[\kappa(\theta) = \frac{2(\rho')^2 - \rho\rho'' + \rho^2}
		{((\rho')^2+\rho^2)^{3/2}}.\]
	\end{enumerate}
\end{problem*}

\begin{soln}
	(a) The curve can be expressed as
	$\alpha(\theta) = (\rho(\theta)\cos\theta, \rho(\theta)\sin\theta)$.
	Then, its arc length is
	\[\int_a^b |\alpha'(\theta)|\dif\theta.\]
	We see that
	\[\alpha'(\theta) = (-\rho\sin\theta + \rho'\cos\theta,
	\rho\cos\theta + \rho'\sin\theta),\]
	so that we can calculate $|\alpha'(\theta)| = \sqrt{\rho^2 + (\rho')^2}$,
	justifying the given arc length formula.
	
	(b) We use the following formula
	(proof copied over from the previous problem):
	\begin{claim*}[Problem 1.5.12(d)]
		If $\alpha\colon I\to\RR^2$ is a plane curve
		$\alpha(t) = (x(t), y(t))$, the signed curvature of $\alpha$ at $t$ is
		\[\kappa(t) = \frac{x'y'' - x''y'}{((x')^2 + (y')^2)^{3/2}}.\]
	\end{claim*}
	\begin{proof}
		To start with we can calculate the unit tangent
		\[T(t) = \left(\frac{x'}{\sqrt{(x')^2 + (y')^2}},
		\frac{y'}{\sqrt{(x')^2 + (y')^2}}\right).\]
		Then, the derivative of the unit tangent with respect to $t$ is
		\[\left(\frac{-y'(x'y'' - x''y')}{((x')^2 + (y')^2)^{3/2}},
		\frac{x'(x'y'' - x''y')}{((x')^2 + (y')^2)^{3/2}}\right),\]
		which can be verified by some menial computation.
		To get the derivative with respect to $s$, we just multiply this by
		$dt/ds = \frac1{|\alpha'(t)|}$. This means that,
		if $n$ is the normal vector obtained by a $\frac\pi2$-counterclockwise
		rotation of $T$, then the signed curvature is the scalar $\kappa$
		such that $T = \kappa n$: in this case
		\[n = \left(\frac{-y'}{|\alpha'|}, \frac{x'}{|\alpha'|}\right),\]
		and we know that
		\[\frac{dT}{ds} = \frac{1}{|\alpha'|}\cdot
		\frac{x'y'' - x''y'}{((x')^2 + (y')^2)^{3/2}}\cdot(-y',x') =
		\frac{x'y'' - x''y'}{((x')^2 + (y')^2)^{3/2}}\cdot
		\frac{(-y', x')}{|\alpha'|},\]
		so we can see that the curvature is indeed
		the expression in the claim.
	\end{proof}

	We have already compute $x'$ and $y'$ above; we can also calculate
	\begin{align*}
		x'' &= -\rho\cos\theta - \rho'\sin\theta - \rho'\sin\theta
		+ \rho''\cos\theta = -\rho\cos\theta - 2\rho'\sin\theta
		+ \rho''\cos\theta\\
		y'' &= -\rho\sin\theta + \rho'\cos\theta + \rho'\cos\theta
		+ \rho''\sin\theta = -\rho\sin\theta + 2\rho'\cos\theta
		+ \rho''\sin\theta.
	\end{align*}
	Then
	\begin{align*}
		x'y'' &= (-\rho\sin\theta + \rho'\cos\theta)
		(-\rho\sin\theta + 2\rho'\cos\theta + \rho''\sin\theta)\\
		&= \rho^2\sin^2\theta - 3\rho\rho'\cos\theta\sin\theta
		- \rho\rho''\sin^2\theta + 2(\rho')^2\cos^2\theta
		+ \rho'\rho''\cos\theta\sin\theta\\
		y'x'' &= (\rho\cos\theta + \rho'\sin\theta)
		(-\rho\cos\theta - 2\rho'\sin\theta + \rho''\cos\theta)\\
		&= -\rho^2\cos^2\theta - 3\rho\rho'\cos\theta\sin\theta
		+\rho\rho''\cos^2\theta - 2(\rho')^2\sin^2\theta
		+ \rho'\rho''\cos\theta\sin\theta
	\end{align*}
	so the difference is
	\[x'y'' - x''y' = \rho^2 - \rho\rho'' + 2(\rho')^2,\]
	hence the signed curvature is
	\[\frac{x'y'' - x''y'}{|\alpha'|^3} =
	\frac{\rho^2 - \rho\rho'' + 2(\rho')^2}{(\rho^2 + (\rho')^2)^{3/2}},\]
	but this is equal to the unsigned curvature
	because the numerator is always nonnegative, for the expression
	$a^2 - ab + 2b^2$ is always nonnegative $a,b$ as it's equal to
	a linear combination of squares:
	\[\left(a - \frac b2\right)^2 + \frac{7b^2}{4}.\qedhere\]
\end{soln}

\begin{problem*}1.7.1.
	Is there a simple closed curve in the plane
	with length equal to $6$ feet and bounding an area of $3$ square feet?
\end{problem*}

\begin{soln}
	\framebox{No}; the isoperimetric inequality says that
	\[\ell^2 \ge 4\pi A\]
	must hold, where $\ell$ is the length of the curve and $A$ is
	the area bounded by the curve, but we can see clearly that
	\[6^2 = 12\cdot 3 < 4\pi\cdot 3 = 12\pi,\]
	so this combination of area and arc length is impossible.
\end{soln}

\begin{problem*}1.7.2.
	Let $\ol{AB}$ be a segment of a straight line and let $\ell$
	be greater than the length of $\ol{AB}$. The curve $\gamma$
	with endpoints $A$ and $B$ has length $\ell$ and maximizes
	the area bounded by $\gamma$ and segment $\ol{AB}$.
	Show that $\gamma$ is an arc of a circle passing through $A$ and $B$.
\end{problem*}

\begin{soln}
	We will take for granted that $\gamma\colon[0,1]\to\RR^2$
	with $\gamma(0) = A$, $\gamma(1) = B$ is a
	parametrized non self-intersecting regular curve.
	In an abuse of notation, we will also refer to the image of $\gamma$
	by $\gamma$ (so for example when we say a point $P$ ``lies on $\gamma$''
	we mean that there is a $t\in[0,1]$ such that $\gamma(t) = P$).
	
	Also, let $\gamma'\colon[0,2]\to\RR^2$ be the curve such that
	$\gamma'|_{[0,1]} = \gamma$ and $\gamma'|_{[1,2]}$ is a parametrization
	of segment $\ol{AB}$ with $\gamma'(1) = B$, $\gamma'(2) = \gamma'(0) = A$.
	Then $\gamma'$ is a simple closed piecewise regular parametrized curve.
	Let $R$ be the region whose area we are trying to maximize
	(i.e. the region bounded by $\gamma'$).
	We first show the following claim:
	\begin{claim*}
		If line $\ol{AB}$ intersects the interior of $R$, then
		the area of $R$ is not maximal.
	\end{claim*}
	\begin{proof}
		Suppose line $\ol{AB}$ intersects the interior of $R$.
		and orient the plane so that $\ol{AB}$ lies in the $x$-axis
		with $A$ to the left of $B$ (so as $t$ increases,
		the point $\gamma'(t)$ moves ``clockwise'' along the curve).
		Moreover, $\ol{AB}$ borders an ``inside region'' and an
		``outside region''; if necessary, reflect the plane about the $x$-axis
		so that the side of $\ol{AB}$ bordering the ``outside region''
		is below the $x$-axis. That is, if an ant were to stand at
		an arbitrary point on segment $\ol{AB}$ (but not on $A$ or $B$)
		and then walk $\eps$ units in the negative $y$ direction
		for any sufficiently small $\eps > 0$, it would be outside $R$.
		
		Let $T(t)$ be the tangent line to $\gamma$ at the point $\gamma(t)$.
		As $t$ varies from $0$ to $1$, this line varies continuously.
		If $T(0)$ does not intersect the interior of $R$, then
		one can rotate this line counterclockwise until it first intersects
		$\gamma$ at a point lying below the $x$-axis;
		let $\lambda$ be the resulting line.
		
		Otherwise, if $T(0)$ does intersect the interior of $R$, then
		consider the infimum $t_0$ of all $t$ such that $T(t)$
		does not intersect the interior of $R$ (there exist such $t$
		because the tangent line at the point of $\gamma$
		with maximal $y$-coordinate is necessarily horizontal
		and thus can't intersect the interior of $R$).
		The line $T(t_0)$ intersects $\gamma$ at $P = \gamma(t_0)$ and
		another point $Q$ such that $P$, $A$, $B$, and $Q$
		lie on $\gamma'$ in counterclockwise order.
		In this case let $\lambda = T(t_0)$.
		
		In either case, we can reflect $\ol{AB}$
		and the arcs joining $A$ and $B$ to the nearest points of $\gamma$
		lying on $\lambda$ over the line $\lambda$.
		Since $\gamma$ can only intersect the $x$-axis finitely many times
		except at arcs lying entirely within the $x$-axis
		(maybe more accurately phrased as ``can only start intersecting
		the $x$-axis finitely many times''), we can repeat this process
		until eventually the resulting curve is entirely within
		one of the half-planes generated by the image of line $\ol{AB}$,
		and each reflection increases the area of the bounded region
		while preserving the arc length; thus we have constructed
		a $\widehat\gamma$ that bounds a larger area than that of $R$.
	\end{proof}
	
	Now let $O$ be a point such that $OA = OB = r$
	and $O$ and $C$ are on different sides of line $AB$.
	If $2\theta = \angle AOB$, then we impose the restriction
	\[\frac{2\theta}{2\pi} = \frac{\ell}{2\pi r}\iff
	\ell = 2r\theta.\]
	To see that such an $(r, \theta)$ pair exists, we can just
	compute directly that
	\[AB = 2r\sin\theta = \frac{\ell\sin\theta}{\theta}.\]
	This expression achieves $0$ at $\theta = \pi/2$ and is equal to $\ell$
	as $\theta$ approaches $0$, meaning $\frac{\ell\sin\theta}{\theta}$
	achieves every value between $[0,\ell)$ as $\theta$ varies
	over $[0,\pi/2]$; the fact that $AB < \ell$ then
	guarantees the existence of a solution by the intermediate value theorem.
	
	This means that if $C'$ is the arc centered at $O$
	with endpoints $A$ and $B$ lying on the opposite side of $AB$ as $C$,
	then the closed curve $\gamma$ composed of $C$ and $C'$
	is a piecewise regular curve with fixed arc length $2\pi r$.
	Since the area of the region bounded by $C'$ and segment $AB$ is fixed,
	maximizing the area of $R$ is equivalent to maximizing the area
	of the region bounded by $\gamma$. By the isoperimetric inequality,
	the area is maximized when $\gamma$ is the circle centered at $O$
	with radius $r$. which exactly implies that $C$ is an arc of a circle
	passing through $A$ and $B$ (an arc of $\gamma$ to be precise).
\end{soln}

(what is the nice way to do this problem? it caused me too much pain)























\end{document}